\makeatletter
\newcommand{\DSTUAsbuk}[1]{\expandafter\@DSTUAsbuk\csname c@#1\endcsname}
\def\@DSTUAsbuk#1{\ifcase#1\or А\or Б\or В\or Г\or Д\or Е\or Ж\or И\or К\or Л\or М\or Н\or П\or Р\or С\or Т\or У\or Ф\or Х\or Ц\or Ш\or Щ\or Ю\or Я\fi}
\makeatother

\appendix
\clearpage

% Auto Cyrillic numbering per ДСТУ 3008:2015 (skips Ґ, Є, З, І, Ї, Й, О, Ч, Ь)
\renewcommand{\thesection}{\DSTUAsbuk{section}}
\titleformat{\section}
  {\clearpage\centering\bfseries\Large}
  {Додаток~\thesection.}
  {0.5em}
  {}

\renewcommand{\thesubsection}{\thesection.\arabic{subsection}}
\titleformat{\subsection}
  {\bfseries\normalsize}
  {\thesubsection.}
  {0.5em}
  {}

% Show only section-level entries in ToC for appendices
\addtocontents{toc}{\protect\setcounter{tocdepth}{1}}
\addtocontents{toc}{\protect\renewcommand*{\protect\cftsecpresnum}{Додаток~}}
\addtocontents{toc}{\protect\renewcommand*{\protect\cftsecaftersnum}{.}}
\addtocontents{toc}{\protect\setlength{\protect\cftsecnumwidth}{5.5em}}

\section*{\centering ДОДАТКИ}\addcontentsline{toc}{section}{ДОДАТКИ}

% ============================================================
% Додаток А: Обов'язковий додаток (§13 Наказу МОН №40)
% ============================================================
\section{Список публікацій здобувача за темою дисертації та відомості про апробацію результатів дисертації}\label{app:publications}

\printpratsi

\setlength \tabcolsep {3pt}

\begin{center}
\begin{spacing}{1}

\begin{longtable}{|p{0.5\linewidth}|p{0.45\linewidth}|}
  \caption{Відомості про апробацію результатів дисертації}
  \label{tab:approbation} \\

\hline
\hfil \textbf{Назва заходу} &

\hfil \begin{tabular}{c}
     \textbf{Місце та дата}  \\
     \textbf{проведення}
\end{tabular}

\\

\hline

\endfirsthead

\multicolumn{2}{r}%
{{\itshape Продовження таблиці \thetable{} }} \\
\hline

\hfil \textbf{Назва заходу} &

\hfil \begin{tabular}{c}
     \textbf{Місце та дата}  \\
     \textbf{проведення}
\end{tabular}

\\

\hline

\endhead

\endfoot

\endlastfoot

IV Міжнародна науково-практична конференція <<Topical Issues of Modern Science, Society and Education>> & Харків, 1--3~листопада 2021~р. \\\hline

VІІ Всеукраїнська науково-практична конференція <<Формула творчості: теорія і методика мистецької освіти>> & Полтава, 26~квітня 2023~р. \\\hline

Всеукраїнська науково-практична конференція <<Актуальні проблеми формування естетичної культури учнівської та студентської молоді в закладах освіти України>> & Кривий Ріг, 27~квітня 2023~р. \\\hline

І~Міжнародна науково-практична конференція Таврійського національного університету до 160-ї річниці від дня народження В.\,І.~Вернадського & Київ, 16--17~березня 2023~р. \\\hline

Щорічна науково-практична конференція <<Образотворче мистецтво як інструмент міжкультурної комунікації в умовах глобалізації>> & Кривий Ріг, 24~лютого 2026~р. \\\hline

\end{longtable}

\end{spacing}
\end{center}

% ============================================================
% Додаток Б: Інструменти опитування та дані
% ============================================================
\section{Інструменти опитування та таблиці первинних даних}\label{app:surveys}

\subsection{Демографічна характеристика вибірки}

Опитування проведено у березні 2024~р. на платформі MOODLE Криворізького державного педагогічного університету. Загальна кількість респондентів представлена у таблиці~\ref{tab:app-demographics}.

\begin{table}[htbp]
\centering
\caption{Демографічна характеристика вибірки опитувань}\label{tab:app-demographics}
\small
\begin{tabularx}{\textwidth}{@{}Xrrr@{}}
\toprule
\textbf{Показник} & \textbf{Випускники} & \textbf{Роботодавці} & \textbf{Студенти (Moodle)} \\
\midrule
Загальна кількість & 32 & 9 & 61 \\
\addlinespace
\textit{За програмами:} & & & \\
\quad Графічний дизайн (ГД) & 12 & 5 & --- \\
\quad Дизайн одягу (ДО) & 13 & 2 & --- \\
\quad Дизайн середовища (ДС) & 7 & 2 & --- \\
\addlinespace
Кількість програм & 3 & 3 & 3 \\
Кількість запитань (Moodle) & --- & --- & 50 \\
\bottomrule
\end{tabularx}
\end{table}

\subsection{Опитувальник випускників: якість підготовки}

Випускникам пропонувалося оцінити рівень підготовки за освітньою програмою за такими вимірами якості (шкала від 1~--- <<не задоволені>> до 5~--- <<дуже задоволені>>):

\begin{enumerate}
    \item Загальнотеоретична підготовка.
    \item Спеціальна (професійна) підготовка.
    \item Практична підготовка.
    \item Формування комунікативної компетентності.
    \item Розвиток здатності до навчання й саморозвитку.
    \item Формування умінь працювати в колективі, команді.
    \item Розвиток ерудованості, загальної культури.
    \item Оволодіння інформаційно-комунікаційними технологіями.
    \item Розвиток лідерських якостей.
\end{enumerate}

Додатково випускники оцінювали актуальність освітньої програми за 5-бальною шкалою та надавали відкриті пропозиції щодо покращення підготовки.

\subsection{Опитувальник роботодавців: оцінка якості випускників}

Роботодавці оцінювали якість та навички випускників за 5-бальною шкалою (1~--- низький рівень, 5~--- високий рівень) за такими критеріями:

\begin{enumerate}
    \item Загальнотеоретична підготовка.
    \item Спеціальна (професійна) підготовка.
    \item Комунікативна компетентність.
    \item Здатність до навчання й саморозвитку.
    \item Адаптація до нових умов.
    \item Стратегічне мислення.
    \item Здатність працювати в колективі, команді.
    \item Здатність ефективно представляти себе й результати своєї праці.
    \item Націленість на професійний розвиток і кар'єрне зростання.
    \item Ерудованість, загальна культура.
    \item Володіння інформаційно-комунікаційними технологіями.
    \item Вміння логічно мислити, робити правильні висновки.
    \item Знання іноземних мов.
\end{enumerate}

Додатково роботодавці визначали пріоритетність напрямків покращення підготовки (шкала 1--5): покращення матеріально-технічної бази, підвищення кваліфікації викладачів, збільшення обсягу практичної підготовки, зміни змісту освітніх програм, розширення співпраці з роботодавцями.

\subsection{Опитування студентів на платформі Moodle}

Анонімне опитування на платформі Moodle КДПУ охоплювало 50 запитань за такими тематичними блоками, які оцінювалися за 5-бальною шкалою Лікерта:

\begin{enumerate}
    \item Відповідність навчання очікуванням (питання 1).
    \item Задоволеність рівнем знань та умінь (питання 3).
    \item Дієвість індивідуальної освітньої траєкторії (питання 7).
    \item Організація та зміст практик (питання 12).
    \item Форми й методи навчання та викладання (питання 19).
    \item Студентоцентрованість викладання (питання 20).
    \item Прозорість контрольних заходів (питання 28).
    \item Професіоналізм викладачів (питання 34).
    \item Матеріально-технічне забезпечення (питання 37).
    \item Інформаційне та навчально-методичне забезпечення (питання 39).
    \item Актуальність освітньої програми (питання 49).
\end{enumerate}

\subsection{Результати тесту Крускала--Волліса}

Для визначення статистично значущих відмінностей між трьома програмами (ГД, ДО, ДС) за результатами Moodle-опитування застосовано непараметричний тест Крускала--Волліса ($H$-статистика, $\alpha = 0{,}05$). Загальна кількість респондентів $n = 61$.

\begin{longtable}{@{}p{0.04\textwidth}p{0.49\textwidth}rrp{0.08\textwidth}@{}}
\caption{Результати тесту Крускала--Волліса для порівняння програм ($n=61$, $k=3$)}\label{tab:app-kruskal} \\
\hline
\textbf{№} & \textbf{Запитання} & \textbf{$H$} & \textbf{$p$} & \textbf{Знач.} \\
\hline
\endfirsthead
\hline
\textbf{№} & \textbf{Запитання} & \textbf{$H$} & \textbf{$p$} & \textbf{Знач.} \\
\hline
\endhead
1 & Чи відповідає навчання на освітній програмі Вашим очікуванням? & 0,407 & 0,816 & Ні \\
\addlinespace
3 & Чи задоволені Ви рівнем отримуваних знань та умінь? & 5,495 & 0,064 & Ні \\
\addlinespace
7 & Як Ви оцінюєте дієвість процедур формування індивідуальної освітньої траєкторії? & 3,347 & 0,188 & Ні \\
\addlinespace
12 & Чи задовольняє Вас рівень організації та зміст практик? & 2,344 & 0,310 & Ні \\
\addlinespace
19 & Чи задоволені Ви формами й методами навчання та викладання на освітній програмі? & 0,682 & 0,711 & Ні \\
\addlinespace
20 & Чи погоджуєтеся Ви з твердженням, що викладання в Університеті є студентоцентрованим? & 3,720 & 0,156 & Ні \\
\addlinespace
28 & Чи вважаєте Ви контрольні заходи прозорими і зрозумілими? & 1,033 & 0,597 & Ні \\
\addlinespace
34 & Чи вважаєте Ви викладачів, які у Вас викладають, професіоналами високого рівня? & 4,218 & 0,121 & Ні \\
\addlinespace
\textbf{37} & \textbf{Чи задоволені Ви рівнем матеріально-технічного забезпечення навчання на освітній програмі?} & \textbf{8,708} & \textbf{0,013} & \textbf{Так} \\
\addlinespace
\textbf{39} & \textbf{Чи задоволені Ви рівнем інформаційного та навчально-методичного забезпечення навчання?} & \textbf{6,019} & \textbf{0,049} & \textbf{Так} \\
\addlinespace
49 & Чи вважаєте Ви освітню програму, за якою навчаєтеся, актуальною та затребуваною? & 0,306 & 0,858 & Ні \\
\hline
\end{longtable}

\textit{Примітка.} Статистично значущими ($p < 0{,}05$) виявилися відмінності за питаннями 37 (матеріально-технічне забезпечення, $p = 0{,}013$) та 39 (інформаційне забезпечення, $p = 0{,}049$). За всіма іншими 9 питаннями відмінності між програмами є несуттєвими, що свідчить про подібний рівень задоволеності студентів ГД, ДО та ДС більшістю аспектів навчання.

\clearpage

% ============================================================
% Додаток В: Аналіз НАЗЯВО
% ============================================================
\section{Повні дані аналізу НАЗЯВО}\label{app:naqa-data}

\subsection{Зведена характеристика 122 програм спеціальності 022 Дизайн}

\begin{table}[htbp]
\centering
\caption{Розподіл освітніх програм спеціальності 022 Дизайн за рівнями}\label{tab:app-naqa-levels}
\small
\begin{tabularx}{\textwidth}{@{}Xrrrr@{}}
\toprule
\textbf{Показник} & \textbf{Бакалавр} & \textbf{Магістр} & \textbf{PhD} & \textbf{Усього} \\
\midrule
Кількість програм & 74 & 43 & 5 & 122 \\
Частка від загальної кількості, \% & 60,7 & 35,2 & 4,1 & 100,0 \\
Кількість компонентів & 2\,442 & 559 & 42 & 3\,043 \\
Середня кількість дисциплін на програму & 33,0 & 13,0 & 8,8 & 24,9 \\
Стандартне відхилення & 12,1 & 6,4 & 3,2 & 12,9 \\
Діапазон дисциплін & 12--73 & 7--38 & 5--14 & 7--73 \\
Медіана & 32 & 12 & 8 & 24 \\
Частка вибіркових компонентів & $\approx$0\% & $\approx$0\% & $\approx$0\% & $\approx$0\% \\
Покриття силабусами & 99,0\% & 98,5\% & 97,6\% & 99,0\% \\
\bottomrule
\end{tabularx}
\end{table}

\begin{table}[htbp]
\centering
\caption{Структура освітніх компонентів 122 програм}\label{tab:app-naqa-components}
\small
\begin{tabularx}{\textwidth}{@{}Xrrr@{}}
\toprule
\textbf{Тип компонента} & \textbf{Кількість} & \textbf{Частка, \%} & \textbf{Унікальних назв} \\
\midrule
Обов'язкові дисципліни & 2\,602 & 85,5 & --- \\
Практики & 335 & 11,0 & --- \\
Атестації & 105 & 3,5 & --- \\
Вибіркові & 1 & $<$0,1 & --- \\
\addlinespace
\textbf{Усього} & \textbf{3\,043} & \textbf{100,0} & \textbf{1\,515} \\
\bottomrule
\end{tabularx}
\end{table}

\subsection{Матриця аналізу ШІ/ІКК-розривів}

Нижче наведено повну матрицю покриття ШІ-навичок та ІКК-компонентів у 122 програмах (за результатами класифікатора на основі 20 ключових слів).

\begin{table}[htbp]
\centering
\caption{Матриця аналізу ШІ/ІКК-розривів у 122 програмах дизайну}\label{tab:app-gap-matrix}
\small
\begin{tabularx}{\textwidth}{@{}Xccl@{}}
\toprule
\textbf{Навичка / компонент} & \textbf{Покриття, \%} & \textbf{Статус} & \textbf{Коментар} \\
\midrule
\multicolumn{4}{@{}l}{\textit{Цифрові інструменти:}} \\
\quad 3D-моделювання & 100 & \checkmark & Присутнє в усіх кластерах \\
\quad Моушн-дизайн / Анімація & 100 & \checkmark & Стабільне покриття \\
\quad UX/UI дизайн & 100 & \checkmark & Зростаюча тенденція \\
\quad Інтерактивний дизайн & 100 & \checkmark & Обчислювальний підхід \\
\quad 2D цифрові інструменти & 33,3 & Частково & Лише Adobe Suite \\
\quad Візуалізація даних & 33,3 & Частково & Окремі програми \\
\quad VR/AR дизайн & 33,3 & Частково & Експериментально \\
\addlinespace
\multicolumn{4}{@{}l}{\textit{Компоненти ШІ:}} \\
\quad Основи ШІ & \textbf{0} & \ding{55} & Повна відсутність \\
\quad Генеративний ШІ та дизайн & \textbf{0} & \ding{55} & Повна відсутність \\
\quad Комп'ютерний зір & \textbf{0} & \ding{55} & Повна відсутність \\
\addlinespace
\midrule
\textbf{Покриття цифрових інстр.} & \textbf{71,4} & Задовільно & 5 із 7 \\
\textbf{Покриття ШІ-компонентів} & \textbf{0,0} & Критично & 0 із 3 \\
\textbf{Загальне покриття} & \textbf{35,7} & Незадовільно & 5 із 14 \\
\bottomrule
\end{tabularx}
\end{table}

\subsection{Метрики мережевого аналізу}

\begin{table}[htbp]
\centering
\caption{Характеристики мережі спільної зустрічальності дисциплін}\label{tab:app-network-metrics}
\small
\begin{tabularx}{\textwidth}{@{}Xr@{}}
\toprule
\textbf{Метрика} & \textbf{Значення} \\
\midrule
Кількість вузлів (унікальних дисциплін) & 1\,240 \\
Кількість ребер (зв'язків) & 7\,524 \\
Щільність мережі ($D$) & 0,0098 \\
Коефіцієнт кластеризації ($C$) & 0,3641 \\
Кількість зв'язаних компонентів & 732 \\
Найбільший зв'язаний компонент (вузлів) & 507 \\
Кількість спільнот (Louvain) & 10 \\
\bottomrule
\end{tabularx}
\end{table}

\begin{table}[htbp]
\centering
\caption{10 спільнот у мережі навчальних програм (Louvain)}\label{tab:app-communities}
\small
\begin{tabularx}{\textwidth}{@{}rrX@{}}
\toprule
\textbf{№} & \textbf{Обсяг} & \textbf{Ядро спільноти (центральні дисципліни)} \\
\midrule
1 & 69 & Практика, Графіка, Проєктна робота~--- \textit{прикладне ядро} \\
3 & 56 & Живопис, Рисунок, Дизайн/Проєктування~--- \textit{традиційне мистецьке ядро} \\
6 & 49 & Наукові дослідження, Теорія дизайну, Ергономіка~--- \textit{теоретичний кластер} \\
7 & 49 & Філософія, Іноземна мова, Українська мова~--- \textit{загальноосвітній кластер} \\
\bottomrule
\end{tabularx}
\end{table}

\begin{table}[htbp]
\centering
\caption{Найбільш центральні дисципліни за ступенем (degree centrality)}\label{tab:app-centrality}
\small
\begin{tabularx}{\textwidth}{@{}Xrl@{}}
\toprule
\textbf{Дисципліна} & \textbf{Ступінь} & \textbf{Коментар} \\
\midrule
Філософія & 315 & Обов'язковий загальноосвітній компонент \\
Живопис & 247 & Класичне мистецьке ядро \\
Рисунок & 225 & Класичне мистецьке ядро \\
Практика & 204 & Обов'язковий практичний компонент \\
\bottomrule
\end{tabularx}
\end{table}

\textit{Примітка.} Жодна зі спільнот не має центру у цифрових або ШІ-дисциплінах. Цифрова/ШІ-тематика не утворила самостійного кластера в українських програмах дизайну і залишається периферійною.

\clearpage

% ============================================================
% Додаток Г: Навчальний план 240 кредитів
% ============================================================
\section{Повний навчальний план на 240 кредитів ЄКТС}\label{app:curriculum}

Нижче представлено інтегровану модель 4-річного (240~кредитів ЄКТС) навчального плану з наскрізною ШІ-інтеграцією для спеціальності 022~Дизайн. Для кожної дисципліни визначено рівень ШІ-інтеграції: <<Баз.>>~--- базовий, <<Інстр.>>~--- інструментальний, <<Транс.>>~--- трансформаційний. Нові ШІ-орієнтовані дисципліни виділено \textbf{напівжирним шрифтом}.

\begin{longtable}{@{}p{0.03\textwidth}p{0.32\textwidth}p{0.05\textwidth}p{0.07\textwidth}p{0.40\textwidth}@{}}
\caption{Навчальний план: Семестр~1 (30~кредитів)}\label{tab:app-sem1} \\
\hline
\textbf{№} & \textbf{Дисципліна} & \textbf{Кр.} & \textbf{Рівень ШІ} & \textbf{ШІ-компонент} \\
\hline
\endfirsthead
\hline
\textbf{№} & \textbf{Дисципліна} & \textbf{Кр.} & \textbf{Рівень ШІ} & \textbf{ШІ-компонент} \\
\hline
\endhead
1 & Основи візуальної комунікації & 6 & --- & Традиційний курс \\
2 & Рисунок і живопис & 6 & Баз. & ШІ-референси; порівняння <<ручне vs. ШІ>> \\
3 & Історія мистецтва та дизайну & 4 & Баз. & ШІ <<у стилі>>; критичний аналіз \\
4 & \textbf{Цифрова грамотність та основи AI} & \textbf{4} & \textbf{Транс.} & \textbf{Основи ШІ; промпт-інженерія; етика} \\
5 & Типографіка & 5 & Баз. & ШІ-генерація шрифтових пар \\
6 & Іноземна мова & 3 & Баз. & ШІ-термінологія англійською \\
\addlinespace
& \textit{Підсумок семестру~1:} & \textit{28} & & \\
\hline
\end{longtable}

\begin{longtable}{@{}p{0.03\textwidth}p{0.32\textwidth}p{0.05\textwidth}p{0.07\textwidth}p{0.40\textwidth}@{}}
\caption{Навчальний план: Семестр~2 (32~кредити)}\label{tab:app-sem2} \\
\hline
\textbf{№} & \textbf{Дисципліна} & \textbf{Кр.} & \textbf{Рівень ШІ} & \textbf{ШІ-компонент} \\
\hline
\endfirsthead
\hline
\textbf{№} & \textbf{Дисципліна} & \textbf{Кр.} & \textbf{Рівень ШІ} & \textbf{ШІ-компонент} \\
\hline
\endhead
7 & Колористика & 5 & Баз. & ШІ-палітри; кольоровий bias \\
8 & Комп'ютерна графіка (Adobe Suite) & 6 & Інстр. & Adobe Firefly; Generative Fill \\
9 & Дизайн-мислення & 4 & Баз. & ШІ на етапі ідеації \\
10 & Основи композиції & 5 & Інстр. & ШІ-генерація варіантів композиції \\
11 & Академічне письмо & 3 & Баз. & ШІ як асистент; етика \\
12 & Вибіркова 1 & 4 & --- & --- \\
\addlinespace
& \textit{Підсумок семестру~2:} & \textit{27} & & \\
& \textit{Підсумок року~1: 55~кр. + вибіркові = 60~кр.} & & & \\
\hline
\end{longtable}

\begin{longtable}{@{}p{0.03\textwidth}p{0.32\textwidth}p{0.05\textwidth}p{0.07\textwidth}p{0.40\textwidth}@{}}
\caption{Навчальний план: Семестр~3 (30~кредитів)}\label{tab:app-sem3} \\
\hline
\textbf{№} & \textbf{Дисципліна} & \textbf{Кр.} & \textbf{Рівень ШІ} & \textbf{ШІ-компонент} \\
\hline
\endfirsthead
\hline
\textbf{№} & \textbf{Дисципліна} & \textbf{Кр.} & \textbf{Рівень ШІ} & \textbf{ШІ-компонент} \\
\hline
\endhead
13 & Графічний дизайн & 6 & Інстр. & ШІ для логотипів; промпт-інженерія \\
14 & UX/UI дизайн: Основи & 6 & Інстр. & Figma AI; ШІ-wireframes; UX з ШІ \\
15 & 3D-моделювання & 5 & Інстр. & ШІ-генерація 3D; ШІ-текстурування \\
16 & Фотографія для дизайнерів & 4 & Баз. & ШІ-обробка; генеративне редагування \\
17 & Веб-технології & 4 & Баз. & ШІ-асистенти для коду \\
18 & Вибіркова 2 & 5 & --- & --- \\
\addlinespace
& \textit{Підсумок семестру~3:} & \textit{30} & & \\
\hline
\end{longtable}

\begin{longtable}{@{}p{0.03\textwidth}p{0.32\textwidth}p{0.05\textwidth}p{0.07\textwidth}p{0.40\textwidth}@{}}
\caption{Навчальний план: Семестр~4 (30~кредитів)}\label{tab:app-sem4} \\
\hline
\textbf{№} & \textbf{Дисципліна} & \textbf{Кр.} & \textbf{Рівень ШІ} & \textbf{ШІ-компонент} \\
\hline
\endfirsthead
\hline
\textbf{№} & \textbf{Дисципліна} & \textbf{Кр.} & \textbf{Рівень ШІ} & \textbf{ШІ-компонент} \\
\hline
\endhead
19 & Figma та прототипування & 5 & Інстр. & Figma AI; ШІ-плагіни \\
20 & Моушн-дизайн & 5 & Інстр. & RunwayML; ШІ-анімація \\
21 & Ілюстрація & 5 & Інстр. & ШІ-ілюстрація; авторський стиль vs. ШІ \\
22 & Інформаційний дизайн & 4 & Інстр. & ШІ-візуалізація даних \\
23 & Навчальна практика & 6 & Інстр. & Фіксація ШІ на підприємстві \\
24 & Вибіркова 3 & 5 & --- & --- \\
\addlinespace
& \textit{Підсумок семестру~4:} & \textit{30} & & \\
& \textit{Підсумок року~2: 60~кредитів} & & & \\
\hline
\end{longtable}

\begin{longtable}{@{}p{0.03\textwidth}p{0.32\textwidth}p{0.05\textwidth}p{0.07\textwidth}p{0.40\textwidth}@{}}
\caption{Навчальний план: Семестр~5 (30~кредитів)}\label{tab:app-sem5} \\
\hline
\textbf{№} & \textbf{Дисципліна} & \textbf{Кр.} & \textbf{Рівень ШІ} & \textbf{ШІ-компонент} \\
\hline
\endfirsthead
\hline
\textbf{№} & \textbf{Дисципліна} & \textbf{Кр.} & \textbf{Рівень ШІ} & \textbf{ШІ-компонент} \\
\hline
\endhead
25 & \textbf{Генеративний дизайн та AI} & \textbf{6} & \textbf{Транс.} & \textbf{Промпт-інженерія; Stable Diffusion; ControlNet} \\
26 & \textbf{Взаємодія людина--AI в дизайні} & \textbf{4} & \textbf{Транс.} & \textbf{HCI для ШІ; колаборативна творчість} \\
27 & Дизайн бренду: Продвинутий & 5 & Інстр. & ШІ у брендингу; ШІ-стратегії \\
28 & Параметричний та інтерактивний дизайн & 5 & Інстр. & Алгоритмічний дизайн з ШІ \\
29 & Дослідження в дизайні & 4 & Інстр. & ШІ для аналізу даних \\
30 & Вибіркова 4 & 6 & --- & --- \\
\addlinespace
& \textit{Підсумок семестру~5:} & \textit{30} & & \\
\hline
\end{longtable}

\begin{longtable}{@{}p{0.03\textwidth}p{0.32\textwidth}p{0.05\textwidth}p{0.07\textwidth}p{0.40\textwidth}@{}}
\caption{Навчальний план: Семестр~6 (30~кредитів)}\label{tab:app-sem6} \\
\hline
\textbf{№} & \textbf{Дисципліна} & \textbf{Кр.} & \textbf{Рівень ШІ} & \textbf{ШІ-компонент} \\
\hline
\endfirsthead
\hline
\textbf{№} & \textbf{Дисципліна} & \textbf{Кр.} & \textbf{Рівень ШІ} & \textbf{ШІ-компонент} \\
\hline
\endhead
31 & UX/UI: Продвинутий & 5 & Транс. & Повний UX/UI процес з ШІ \\
32 & Візуалізація даних & 5 & Інстр. & ШІ-візуалізація; автоматизація \\
33 & Дизайн середовища & 4 & Баз. & ШІ-візуалізація інтер'єрів \\
34 & Виробнича практика & 9 & Транс. & ШІ на підприємстві; ШІ-портфоліо \\
35 & Вибіркова 5 & 7 & --- & --- \\
\addlinespace
& \textit{Підсумок семестру~6:} & \textit{30} & & \\
& \textit{Підсумок року~3: 60~кредитів} & & & \\
\hline
\end{longtable}

\begin{longtable}{@{}p{0.03\textwidth}p{0.32\textwidth}p{0.05\textwidth}p{0.07\textwidth}p{0.40\textwidth}@{}}
\caption{Навчальний план: Семестр~7 (30~кредитів)}\label{tab:app-sem7} \\
\hline
\textbf{№} & \textbf{Дисципліна} & \textbf{Кр.} & \textbf{Рівень ШІ} & \textbf{ШІ-компонент} \\
\hline
\endfirsthead
\hline
\textbf{№} & \textbf{Дисципліна} & \textbf{Кр.} & \textbf{Рівень ШІ} & \textbf{ШІ-компонент} \\
\hline
\endhead
36 & Дипломне проектування I & 10 & Транс. & ШІ як дослідницький/проєктний інструмент \\
37 & Портфоліо та професійна практика & 4 & Інстр. & Цифрове портфоліо з ШІ-проєктами \\
38 & Дизайн та підприємництво & 4 & Баз. & ШІ-інструменти для бізнесу \\
39 & \textbf{Критичний дизайн та етика AI} & \textbf{4} & \textbf{Транс.} & \textbf{Bias; авторство; deepfake; екологія} \\
40 & Вибіркова спеціалізації 1 & 8 & --- & --- \\
\addlinespace
& \textit{Підсумок семестру~7:} & \textit{30} & & \\
\hline
\end{longtable}

\begin{longtable}{@{}p{0.03\textwidth}p{0.32\textwidth}p{0.05\textwidth}p{0.07\textwidth}p{0.40\textwidth}@{}}
\caption{Навчальний план: Семестр~8 (30~кредитів)}\label{tab:app-sem8} \\
\hline
\textbf{№} & \textbf{Дисципліна} & \textbf{Кр.} & \textbf{Рівень ШІ} & \textbf{ШІ-компонент} \\
\hline
\endfirsthead
\hline
\textbf{№} & \textbf{Дисципліна} & \textbf{Кр.} & \textbf{Рівень ШІ} & \textbf{ШІ-компонент} \\
\hline
\endhead
41 & Дипломне проектування II & 15 & Транс. & ШІ-інтеграція в дипломний проєкт \\
42 & Переддипломна практика & 6 & Транс. & ШІ-стратегія для диплому \\
43 & Вибіркова спеціалізації 2 & 9 & --- & --- \\
\addlinespace
& \textit{Підсумок семестру~8:} & \textit{30} & & \\
& \textit{Підсумок року~4: 60~кредитів} & & & \\
\hline
\end{longtable}

\begin{table}[htbp]
\centering
\caption{Зведена характеристика ШІ-інтеграції у навчальному плані}\label{tab:app-curriculum-summary}
\small
\begin{tabularx}{\textwidth}{@{}Xrrrrr@{}}
\toprule
\textbf{Показник} & \textbf{Рік 1} & \textbf{Рік 2} & \textbf{Рік 3} & \textbf{Рік 4} & \textbf{Усього} \\
\midrule
Загальна кількість дисциплін & 12 & 12 & 11 & 8 & 43 \\
Дисципліни з ШІ-інтеграцією & 10 & 10 & 9 & 6 & 35 \\
\addlinespace
Базовий рівень ШІ & 7 & 2 & 1 & 1 & 11 \\
Інструментальний рівень & 2 & 8 & 4 & 1 & 15 \\
Трансформаційний рівень & 1 & 0 & 4 & 4 & 9 \\
\addlinespace
Нові ШІ-дисципліни & 1 & 0 & 2 & 1 & 4 \\
Кредити нових ШІ-дисциплін & 4 & 0 & 10 & 4 & 18 \\
Кредити ЄКТС за рік & 60 & 60 & 60 & 60 & 240 \\
\bottomrule
\end{tabularx}
\end{table}

\clearpage

% ============================================================
% Додаток Д: Таблиці трансформації дисциплін
% ============================================================
\section{Таблиці трансформації дисциплін КДПУ}\label{app:transforms}

Нижче представлено деталізовані таблиці трансформації дисциплін чотирьох основних груп: студійних, цифрових, теоретичних та професійних курсів. Для кожної дисципліни зазначено поточний стан (до трансформації), рекомендовані зміни (після трансформації), рівень ШІ-інтеграції та цільові компоненти ІКК.

\subsection{Трансформація студійних курсів}

\begin{longtable}{@{}p{0.13\textwidth}p{0.25\textwidth}p{0.30\textwidth}p{0.07\textwidth}p{0.15\textwidth}@{}}
\caption{Деталізована трансформація студійних курсів}\label{tab:app-studio} \\
\hline
\textbf{Дисципліна} & \textbf{До трансформації} & \textbf{Після трансформації} & \textbf{Рівень} & \textbf{Компон. ІКК} \\
\hline
\endfirsthead
\hline
\textbf{Дисципліна} & \textbf{До трансформації} & \textbf{Після трансформації} & \textbf{Рівень} & \textbf{Компон. ІКК} \\
\hline
\endhead
Рисунок і живопис & Традиційна техніка; олівець, акварель, олія; натюрморт, пленер; відсутність цифрового компоненту & + Порівняння ручних робіт з ШІ-генерацією (Midjourney <<in the style of\ldots>>); рефлексія щодо автентичності; ШІ як інструмент пошуку референсів & Баз. & Рефлекс., Інф.-ан. \\
\addlinespace
Основи композиції & Формальна композиція; ритм, контраст, домінанта; виключно ручне виконання & + Генерація 20+ варіантів композиції через ШІ для порівняльного аналізу; вправи <<ручне vs. ШІ>>; документація процесу у рефлексивному щоденнику & Інстр. & Технол., Інф.-ан., Рефлекс. \\
\addlinespace
Колористика & Теорія кольору; кольорові гармонії; ручне змішування; традиційні колірні кола & + ШІ-генерація кольорових палітр (Adobe Color AI, Coolors AI); аналіз кольорових рішень у ШІ-зображеннях; етика кольорового bias & Баз. & Технол., Рефлекс. \\
\hline
\end{longtable}

\textit{Ключовий метод:} порівняльний аналіз <<ручне vs. ШІ>>~--- студент виконує завдання традиційним способом, потім генерує аналогічне за допомогою ШІ, після чого порівнює результати. Обсяг ручної практики зберігається на рівні $\geq$~90\% від поточного.

\subsection{Трансформація курсів цифрового дизайну}

\begin{longtable}{@{}p{0.13\textwidth}p{0.25\textwidth}p{0.30\textwidth}p{0.07\textwidth}p{0.15\textwidth}@{}}
\caption{Деталізована трансформація курсів цифрового дизайну}\label{tab:app-digital} \\
\hline
\textbf{Дисципліна} & \textbf{До трансформації} & \textbf{Після трансформації} & \textbf{Рівень} & \textbf{Компон. ІКК} \\
\hline
\endfirsthead
\hline
\textbf{Дисципліна} & \textbf{До трансформації} & \textbf{Після трансформації} & \textbf{Рівень} & \textbf{Компон. ІКК} \\
\hline
\endhead
Комп'ютерна графіка (Adobe Suite) & Photoshop, Illustrator, InDesign; ручні операції обробки; застаріле ПЗ (за даними опитувань: <<Photoshop 2003>>) & + Adobe Firefly генерація та редагування; Generative Fill/Expand; ШІ-ретуш; автоматизація рутинних операцій; модуль промпт-інженерії (8~год.) & Інстр. & Технол., Інф.-ан. \\
\addlinespace
3D-моделювання & Blender/3ds Max; полігональне моделювання; текстурування; рендеринг & + ШІ-генерація 3D (Meshy, Luma AI); ШІ-текстурування; порівняння ручного та ШІ-моделювання; обговорення обмежень ШІ & Інстр. & Технол., Рефлекс. \\
\addlinespace
Figma та прототипування & Figma; компоненти; прототипи; варіанти; auto layout & + Figma AI для генерації макетів; ШІ-плагіни (Magician, Relume); порівняння ШІ-прототипів з ручними; UX-дослідження з ШІ & Інстр. & Технол., Комунік. \\
\addlinespace
Моушн-дизайн & After Effects; анімація; motion graphics; кейфреймінг & + RunwayML для генеративного відео; ШІ-анімація; промпт-режисура; етика deepfake у дизайні; порівняння підходів & Транс. & Технол., Рефлекс., Комунік. \\
\hline
\end{longtable}

\textit{Ключовий метод:} ШІ-аугментований робочий процес~--- студент спочатку опановує класичний інструмент, потім вивчає інтеграцію ШІ-функцій як частини професійного процесу.

\subsection{Трансформація теоретичних курсів}

\begin{longtable}{@{}p{0.13\textwidth}p{0.25\textwidth}p{0.30\textwidth}p{0.07\textwidth}p{0.15\textwidth}@{}}
\caption{Деталізована трансформація теоретичних курсів}\label{tab:app-theory} \\
\hline
\textbf{Дисципліна} & \textbf{До трансформації} & \textbf{Після трансформації} & \textbf{Рівень} & \textbf{Компон. ІКК} \\
\hline
\endfirsthead
\hline
\textbf{Дисципліна} & \textbf{До трансформації} & \textbf{Після трансформації} & \textbf{Рівень} & \textbf{Компон. ІКК} \\
\hline
\endhead
Історія мистецтва та дизайну & Хронологічний огляд; стилі та напрямки; візуальний аналіз; лекційна форма & + ШІ-генерація <<у стилі>> історичних напрямків; критичний аналіз точності ШІ-інтерпретацій; дискусія <<Чи може ШІ створити мистецтво?>> & Баз. & Інф.-ан., Рефлекс. \\
\addlinespace
Дизайн-мислення & Design Thinking процес; емпатія, ідеація, прототипування; робота в групах & + ШІ на етапі ідеації (генерація варіантів); ШІ для швидкого прототипування; аналіз впливу ШІ на Design Thinking; етика ШІ у дизайні для людей & Інстр. & Інф.-ан., Технол., Рефлекс. \\
\addlinespace
Дослідження в дизайні & Методологія дизайн-досліджень; якісні та кількісні методи; наукова етика & + ШІ для аналізу даних UX-досліджень; ChatGPT для систематизації якісних даних; критична оцінка ШІ-аналітики & Інстр. & Інф.-ан., Рефлекс. \\
\hline
\end{longtable}

\textit{Ключовий метод:} критичний аналіз ШІ-контенту~--- студенти генерують ШІ-контент, пов'язаний з темою курсу, а потім критично аналізують точність, стереотипність та обмеження ШІ-інтерпретації.

\subsection{Трансформація професійних курсів}

\begin{longtable}{@{}p{0.13\textwidth}p{0.25\textwidth}p{0.30\textwidth}p{0.07\textwidth}p{0.15\textwidth}@{}}
\caption{Деталізована трансформація професійних курсів}\label{tab:app-professional} \\
\hline
\textbf{Дисципліна} & \textbf{До трансформації} & \textbf{Після трансформації} & \textbf{Рівень} & \textbf{Компон. ІКК} \\
\hline
\endfirsthead
\hline
\textbf{Дисципліна} & \textbf{До трансформації} & \textbf{Після трансформації} & \textbf{Рівень} & \textbf{Компон. ІКК} \\
\hline
\endhead
UX/UI: Основи & Wireframing; user flows; usability testing; paper prototyping & + ШІ для генерації wireframes (Figma AI, Uizard); ШІ-аналіз user flows; A/B тестування з ШІ; промпт-інженерія для UI-копірайтингу & Інстр. & Технол., Інф.-ан. \\
\addlinespace
UX/UI: Продвинутий & Складні інтерфейси; дослідницькі методи; дизайн-системи; accessibility & + Повний UX/UI процес з ШІ~--- від дослідження до прототипу; ШІ в UX-аналітиці; дизайн для ШІ-інтерфейсів (conversational UI) & Транс. & Технол., Комунік., Рефлекс. \\
\addlinespace
Графічний дизайн & Фірмовий стиль; плакат; поліграфія; ручні ескізи & + ШІ для генерації логотипів та варіантів; промпт-інженерія для стилістичної консистентності; ШІ у робочому процесі дизайн-студії & Інстр. & Технол., Комунік. \\
\addlinespace
Ілюстрація & Авторська ілюстрація; техніки; стилізація; традиційні медіа & + ШІ-ілюстрація (Midjourney, DALL-E); порівняння авторського та ШІ-стилю; етика авторства; ШІ як інструмент ідеації & Інстр. & Технол., Рефлекс. \\
\addlinespace
Дизайн бренду & Брендинг; позиціонування; візуальна ідентичність; brand book & + ШІ для генерації концепцій бренду; ШІ-аналіз конкурентів; ChatGPT для brand voice; мудборди через ШІ & Інстр. & Технол., Інф.-ан., Комунік. \\
\addlinespace
Інформаційний дизайн & Інфографіка; візуалізація даних; діаграми; схеми & + ШІ для візуалізації даних; ChatGPT для аналізу даних; автоматизація інфографіки; критичний аналіз ШІ-візуалізацій & Інстр. & Технол., Інф.-ан. \\
\hline
\end{longtable}

\textit{Ключовий метод:} ШІ-інтегроване проєктне навчання~--- кожен проєкт включає обов'язковий етап використання ШІ-інструментів з документацією процесу (промпти, ітерації, обґрунтування вибору) у <<ШІ-щоденнику проєкту>>.

\subsection{Трансформація практик та загальноосвітніх курсів}

\begin{longtable}{@{}p{0.13\textwidth}p{0.25\textwidth}p{0.30\textwidth}p{0.07\textwidth}p{0.15\textwidth}@{}}
\caption{Трансформація практик та загальноосвітніх курсів}\label{tab:app-practice-general} \\
\hline
\textbf{Дисципліна} & \textbf{До трансформації} & \textbf{Після трансформації} & \textbf{Рівень} & \textbf{Компон. ІКК} \\
\hline
\endfirsthead
\hline
\textbf{Дисципліна} & \textbf{До трансформації} & \textbf{Після трансформації} & \textbf{Рівень} & \textbf{Компон. ІКК} \\
\hline
\endhead
\multicolumn{5}{@{}l}{\textit{Практики:}} \\
\addlinespace
Навчальна практика (2~курс) & Ознайомча; спостереження за роботою дизайнерів; пасивна участь & + Обов'язкова фіксація використання ШІ на підприємстві; ШІ-проєкт під керівництвом наставника; порівняльний звіт & Інстр. & Технол., Інф.-ан. \\
\addlinespace
Виробнича практика (3~курс) & Самостійна робота на підприємстві; виконання завдань; традиційний звіт & + Вибір баз практики з ШІ-використанням; самостійний ШІ-проєкт; документація ШІ-процесів; ШІ-портфоліо & Транс. & Усі \\
\addlinespace
Переддипломна практика (4~курс) & Збір матеріалу для дипломного проєкту; класична методологія & + ШІ як дослідницький інструмент; ШІ-стратегія для диплому; рефлексивний звіт про роль ШІ у проєкті & Транс. & Усі \\
\addlinespace
\multicolumn{5}{@{}l}{\textit{Загальноосвітні курси:}} \\
\addlinespace
Іноземна мова & Загальна англійська; граматика та лексика; читання & + English for Design: фахова англійська з ШІ-термінологією; робота з документацією ШІ-інструментів; промпт-інженерія англійською & Інстр. & Комунік., Технол. \\
\addlinespace
Академічне письмо & Основи наукового стилю; структура тексту; оформлення & + ШІ як асистент написання (ChatGPT, Grammarly AI); критичний аналіз ШІ-текстів; етика ШІ в академічному письмі & Баз. & Інф.-ан., Рефлекс. \\
\addlinespace
Філософія & Історія філософії; онтологія; епістемологія; лекційна форма & + Модуль <<Філософія ШІ>>: свідомість, авторство, творчість; ШІ та творчість; відповідальність за ШІ-рішення & Баз. & Рефлекс. \\
\addlinespace
Дизайн та підприємництво & Бізнес-планування; маркетинг для дизайнерів; freelance & + ШІ-інструменти для маркетингу; ШІ як конкурентна перевага; бізнес-моделі ШІ-дизайну & Баз. & Комунік., Інф.-ан. \\
\hline
\end{longtable}

\clearpage

% ============================================================
% Додаток Е: Відгуки стейкхолдерів та кодування
% ============================================================
\section{Відгуки стейкхолдерів та витяги з акредитаційних звітів}\label{app:stakeholders}

\subsection{Кодування висновків експертних груп НАЗЯВО}

Для аналізу висновків експертних груп НАЗЯВО щодо програм дизайну КДПУ застосовано метод контент-аналізу з кодуванням за 4 компонентами ІКК. Проаналізовано 4 документи: висновок щодо ОП <<Дизайн одягу>> (скорочений) та 3 повних висновки експертних груп щодо програм ГД, ДО, ДС. Результати кодування представлено у таблиці~\ref{tab:app-expert-coding}.

\begin{table}[htbp]
\centering
\caption{Частотний аналіз ІКК-компонентів у висновках експертних груп НАЗЯВО}\label{tab:app-expert-coding}
\small
\begin{tabularx}{\textwidth}{@{}p{0.27\textwidth}rrrrrr@{}}
\toprule
\textbf{Документ} & \textbf{Обсяг, симв.} & \textbf{Інф.-ан.} & \textbf{Комунік.} & \textbf{Технол.} & \textbf{Рефлекс.} & \textbf{Заг. якість} \\
\midrule
Відповіді ДО (скороч.) & 36\,611 & 12 & 6 & 25 & 1 & 83 \\
ГД висновок ЕГ (повний) & 105\,316 & 114 & 23 & 80 & 12 & 199 \\
ДО висновок ЕГ (повний) & 103\,875 & 101 & 31 & 71 & 11 & 212 \\
ДС висновок ЕГ (повний) & 102\,068 & 108 & 24 & 70 & 12 & 194 \\
\addlinespace
\midrule
\textbf{Усього} & \textbf{347\,870} & \textbf{335} & \textbf{84} & \textbf{246} & \textbf{36} & \textbf{688} \\
\bottomrule
\end{tabularx}
\end{table}

\textit{Позначення:} Інф.-ан.~--- інформаційно-аналітичний компонент; Комунік.~--- комунікативний; Технол.~--- технологічний; Рефлекс.~--- рефлексивний; Заг. якість~--- загальні згадки якості освіти.

Основні спостереження:
\begin{itemize}
    \item \textbf{Інформаційно-аналітичний компонент} має найвищу частоту (101--114 згадувань у повних висновках), що свідчить про увагу експертних груп до інформаційної складової підготовки.
    \item \textbf{Технологічний компонент} представлений стабільно (70--80 згадувань), проте переважно через згадки традиційного обладнання та ПЗ, без ШІ-тематики.
    \item \textbf{Комунікативний компонент} має помірну представленість (23--31 згадування).
    \item \textbf{Рефлексивний компонент} має найменшу представленість (11--12 згадувань у повних висновках), що свідчить про відсутність системи цілеспрямованого формування рефлексивного компоненту ІКК.
\end{itemize}

\subsection{Еволюція ІКК-змісту в освітньо-професійних програмах КДПУ}

На основі контент-аналізу ОПП та силабусів КДПУ за 2017--2021~рр. простежено динаміку ІКК-ключових слів за 4 компонентами (таблиця~\ref{tab:app-opp-evolution}).

\begin{table}[htbp]
\centering
\caption{Еволюція ІКК-змісту в ОПП КДПУ (2017--2021)}\label{tab:app-opp-evolution}
\small
\begin{tabularx}{\textwidth}{@{}Xrrrrr@{}}
\toprule
\textbf{ОПП / рік} & \textbf{Інф.-ан.} & \textbf{Комунік.} & \textbf{Технол.} & \textbf{Рефлекс.} & \textbf{Усього} \\
\midrule
ОПП 2017 (єдина) & 12 & 22 & 56 & 26 & 116 \\
ОПП 2018 (єдина) & 12 & 21 & 61 & 28 & 122 \\
\addlinespace
ГД 2019 & 18 & 11 & 38 & 19 & 86 \\
ДО 2019 & 14 & 10 & 43 & 20 & 87 \\
ДС 2019 & 16 & 12 & 39 & 21 & 88 \\
\addlinespace
ГД 2021 & 16 & 11 & 31 & 17 & 75 \\
ДО 2021 & 14 & 11 & 34 & 16 & 75 \\
ДС 2021 & 14 & 12 & 37 & 19 & 82 \\
\bottomrule
\end{tabularx}
\end{table}

Таблиця демонструє \textbf{нисхідний тренд} загальної кількості ІКК-ключових слів: від 116--122 у 2017--2018~рр. до 75--82 у 2021~р. Це свідчить про розмивання, а не посилення ІКК-спрямованого змісту при реформуванні програм. Особливо помітним є зниження технологічного компоненту (з 56--61 до 31--37) та рефлексивного (з 26--28 до 16--19).

\subsection{Ключові рекомендації роботодавців}

За результатами опитування 9 роботодавців трьох програм виділено такі ключові рекомендації (за пріоритетністю):

\begin{enumerate}
    \item \textbf{Збільшення обсягу практичної підготовки}~--- 100\% роботодавців усіх програм.
    \item \textbf{Розширення співпраці з роботодавцями}~--- 100\% роботодавців.
    \item \textbf{Покращення матеріально-технічної бази}~--- 100\% роботодавців ДО, 80\% ГД.
    \item \textbf{Підвищення кваліфікації викладачів}~--- 60--100\% роботодавців.
    \item \textbf{Зміни змісту освітніх програм}~--- 60\% роботодавців.
    \item \textbf{Використання штучного інтелекту у сфері дизайну}~--- пряма рекомендація роботодавців ДО.
    \item \textbf{Покращення іншомовної підготовки}~--- 100\% роботодавців ГД.
\end{enumerate}

Особливо значущою є пряма рекомендація роботодавців програми <<Дизайн одягу>> щодо <<використання штучного інтелекту у сфері дизайну>>, що є прямим підтвердженням запиту ринку праці на ШІ-компетентності випускників-дизайнерів.

\clearpage

% ============================================================
% Додаток Ж: Анкета експертного оцінювання
% ============================================================
\section{Анкета експертного оцінювання моделі}\label{app:expert-survey}

\begin{center}
\textbf{АНКЕТА ЕКСПЕРТНОГО ОЦІНЮВАННЯ}\\[0.3cm]
\textbf{Структурно-функціональної моделі формування}\\
\textbf{інформаційно-комунікативної компетентності (ІКК)}\\
\textbf{майбутніх дизайнерів в умовах інтеграції штучного інтелекту}\\[0.5cm]
\end{center}

\textbf{Шановний(-а) експерте!}

Просимо Вас оцінити структурно-функціональну модель формування інформаційно-комунікативної компетентності (ІКК) майбутніх дизайнерів, розроблену в рамках дисертаційного дослідження. Модель спрямована на системну трансформацію навчальних програм спеціальності 022~Дизайн з інтеграцією інструментів штучного інтелекту.

Будь ласка, оцініть кожне твердження за 5-бальною шкалою:

\begin{center}
\small
\begin{tabular}{cl}
\toprule
\textbf{Бал} & \textbf{Значення} \\
\midrule
1 & Повністю не погоджуюсь \\
2 & Скоріше не погоджуюсь \\
3 & Частково погоджуюсь \\
4 & Скоріше погоджуюсь \\
5 & Повністю погоджуюсь \\
\bottomrule
\end{tabular}
\end{center}

\vspace{0.3cm}

\textbf{Відомості про експерта}\\[0.2cm]
\begin{tabular}{@{}p{0.45\textwidth}p{0.50\textwidth}@{}}
ПІБ: & \rule{0.50\textwidth}{0.4pt} \\[0.3cm]
Науковий ступінь, вчене звання: & \rule{0.50\textwidth}{0.4pt} \\[0.3cm]
Посада, місце роботи: & \rule{0.50\textwidth}{0.4pt} \\[0.3cm]
Стаж у галузі дизайн-освіти (роки): & \rule{0.50\textwidth}{0.4pt} \\[0.3cm]
Досвід експертизи НАЗЯВО: & \rule{0.50\textwidth}{0.4pt} \\[0.3cm]
\end{tabular}

\vspace{0.5cm}

\textbf{Блок~1. Мотиваційно-цільовий блок моделі}

\begin{longtable}{@{}p{0.04\textwidth}p{0.74\textwidth}p{0.04\textwidth}p{0.04\textwidth}p{0.04\textwidth}p{0.04\textwidth}p{0.04\textwidth}@{}}
\hline
\textbf{№} & \textbf{Твердження} & \textbf{1} & \textbf{2} & \textbf{3} & \textbf{4} & \textbf{5} \\
\hline
\endfirsthead
\hline
\textbf{№} & \textbf{Твердження} & \textbf{1} & \textbf{2} & \textbf{3} & \textbf{4} & \textbf{5} \\
\hline
\endhead
1 & Мета моделі (формування ІКК через системну трансформацію навчальних програм з інтеграцією ШІ) є актуальною та обґрунтованою. & $\square$ & $\square$ & $\square$ & $\square$ & $\square$ \\
\addlinespace
2 & Завдання моделі за 4 компонентами ІКК (інформаційно-аналітичний, комунікативний, технологічний, рефлексивний) є конкретними та досяжними. & $\square$ & $\square$ & $\square$ & $\square$ & $\square$ \\
\addlinespace
3 & Соціальне замовлення (дані опитувань роботодавців, випускників, аналіз програм) достатньо обґрунтовує необхідність трансформації. & $\square$ & $\square$ & $\square$ & $\square$ & $\square$ \\
\addlinespace
4 & Формулювання компонентів ІКК є повним та адекватним для дизайнерської професії. & $\square$ & $\square$ & $\square$ & $\square$ & $\square$ \\
\addlinespace
5 & Модель враховує сучасні вимоги ринку праці до дизайнерів. & $\square$ & $\square$ & $\square$ & $\square$ & $\square$ \\
\hline
\end{longtable}

Ваші зауваження та рекомендації щодо блоку~1:

\rule{\textwidth}{0.4pt}\\[0.3cm]
\rule{\textwidth}{0.4pt}\\[0.5cm]

\textbf{Блок~2. Теоретико-методологічний блок моделі}

\begin{longtable}{@{}p{0.04\textwidth}p{0.74\textwidth}p{0.04\textwidth}p{0.04\textwidth}p{0.04\textwidth}p{0.04\textwidth}p{0.04\textwidth}@{}}
\hline
\textbf{№} & \textbf{Твердження} & \textbf{1} & \textbf{2} & \textbf{3} & \textbf{4} & \textbf{5} \\
\hline
6 & Теоретичні підходи (TPACK, студійна педагогіка, компетентнісний підхід) є обґрунтованим підґрунтям моделі. & $\square$ & $\square$ & $\square$ & $\square$ & $\square$ \\
\addlinespace
7 & Дидактичні принципи (системність, доповнення, прогресивність, практикоорієнтованість, рефлексивність, відповідність ринку) є повними та послідовними. & $\square$ & $\square$ & $\square$ & $\square$ & $\square$ \\
\addlinespace
8 & Принцип <<доповнення, а не заміщення>> (збереження ручних навичок) є необхідним для дизайн-освіти. & $\square$ & $\square$ & $\square$ & $\square$ & $\square$ \\
\addlinespace
9 & Трансформація студійної педагогіки від <<виробника>> до <<виробника + куратора>> є педагогічно обґрунтованою. & $\square$ & $\square$ & $\square$ & $\square$ & $\square$ \\
\hline
\end{longtable}

Ваші зауваження та рекомендації щодо блоку~2:

\rule{\textwidth}{0.4pt}\\[0.3cm]
\rule{\textwidth}{0.4pt}\\[0.5cm]

\textbf{Блок~3. Змістовий блок моделі}

\begin{longtable}{@{}p{0.04\textwidth}p{0.74\textwidth}p{0.04\textwidth}p{0.04\textwidth}p{0.04\textwidth}p{0.04\textwidth}p{0.04\textwidth}@{}}
\hline
\textbf{№} & \textbf{Твердження} & \textbf{1} & \textbf{2} & \textbf{3} & \textbf{4} & \textbf{5} \\
\hline
10 & Відображення 4 компонентів ІКК на навчальні дисципліни є коректним та повним. & $\square$ & $\square$ & $\square$ & $\square$ & $\square$ \\
\addlinespace
11 & Рекомендації щодо трансформації студійних курсів (Живопис, Композиція) є реалістичними та ефективними. & $\square$ & $\square$ & $\square$ & $\square$ & $\square$ \\
\addlinespace
12 & Рекомендації щодо трансформації цифрових та професійних курсів є реалістичними та ефективними. & $\square$ & $\square$ & $\square$ & $\square$ & $\square$ \\
\addlinespace
13 & 240-кредитний навчальний план з <<ШІ-спіною>> є збалансованим та реалізовуваним. & $\square$ & $\square$ & $\square$ & $\square$ & $\square$ \\
\hline
\end{longtable}

Ваші зауваження та рекомендації щодо блоку~3:

\rule{\textwidth}{0.4pt}\\[0.3cm]
\rule{\textwidth}{0.4pt}\\[0.5cm]

\textbf{Блок~4. Організаційно-технологічний блок моделі}

\begin{longtable}{@{}p{0.04\textwidth}p{0.74\textwidth}p{0.04\textwidth}p{0.04\textwidth}p{0.04\textwidth}p{0.04\textwidth}p{0.04\textwidth}@{}}
\hline
\textbf{№} & \textbf{Твердження} & \textbf{1} & \textbf{2} & \textbf{3} & \textbf{4} & \textbf{5} \\
\hline
14 & Набір ШІ-інструментів (Midjourney, DALL-E, Stable Diffusion, Adobe Firefly, Figma AI, RunwayML) є актуальним та достатнім. & $\square$ & $\square$ & $\square$ & $\square$ & $\square$ \\
\addlinespace
15 & Методи навчання (модифікована студія, проєктне навчання з ШІ, критика з ШІ-виміром) є ефективними для формування ІКК. & $\square$ & $\square$ & $\square$ & $\square$ & $\square$ \\
\addlinespace
16 & Поетапна стратегія впровадження (4 етапи) є реалістичною. & $\square$ & $\square$ & $\square$ & $\square$ & $\square$ \\
\addlinespace
17 & Ресурсні вимоги (обладнання, кваліфікація викладачів, методичне забезпечення) є адекватними. & $\square$ & $\square$ & $\square$ & $\square$ & $\square$ \\
\hline
\end{longtable}

Ваші зауваження та рекомендації щодо блоку~4:

\rule{\textwidth}{0.4pt}\\[0.3cm]
\rule{\textwidth}{0.4pt}\\[0.5cm]

\textbf{Блок~5. Діагностично-результативний блок моделі}

\begin{longtable}{@{}p{0.04\textwidth}p{0.74\textwidth}p{0.04\textwidth}p{0.04\textwidth}p{0.04\textwidth}p{0.04\textwidth}p{0.04\textwidth}@{}}
\hline
\textbf{№} & \textbf{Твердження} & \textbf{1} & \textbf{2} & \textbf{3} & \textbf{4} & \textbf{5} \\
\hline
18 & Критерії оцінювання ІКК за 4 компонентами є повними та вимірюваними. & $\square$ & $\square$ & $\square$ & $\square$ & $\square$ \\
\addlinespace
19 & Портфоліо-орієнтоване оцінювання є найбільш адекватним інструментом для оцінки ІКК дизайнерів. & $\square$ & $\square$ & $\square$ & $\square$ & $\square$ \\
\addlinespace
20 & 4-рівнева шкала сформованості ІКК (початковий, базовий, достатній, високий) є діагностичною. & $\square$ & $\square$ & $\square$ & $\square$ & $\square$ \\
\hline
\end{longtable}

Ваші зауваження та рекомендації щодо блоку~5:

\rule{\textwidth}{0.4pt}\\[0.3cm]
\rule{\textwidth}{0.4pt}\\[0.5cm]

\textbf{Блок~6. Загальна оцінка моделі}

\begin{longtable}{@{}p{0.04\textwidth}p{0.74\textwidth}p{0.04\textwidth}p{0.04\textwidth}p{0.04\textwidth}p{0.04\textwidth}p{0.04\textwidth}@{}}
\hline
\textbf{№} & \textbf{Твердження} & \textbf{1} & \textbf{2} & \textbf{3} & \textbf{4} & \textbf{5} \\
\hline
21 & Модель загалом є цілісною та внутрішньо узгодженою. & $\square$ & $\square$ & $\square$ & $\square$ & $\square$ \\
\addlinespace
22 & Модель може бути адаптована для інших програм спеціальності 022~Дизайн в Україні. & $\square$ & $\square$ & $\square$ & $\square$ & $\square$ \\
\addlinespace
23 & Реалізація моделі суттєво підвищить рівень сформованості ІКК випускників. & $\square$ & $\square$ & $\square$ & $\square$ & $\square$ \\
\hline
\end{longtable}

Ваші зауваження та рекомендації щодо моделі загалом:

\rule{\textwidth}{0.4pt}\\[0.3cm]
\rule{\textwidth}{0.4pt}\\[0.3cm]
\rule{\textwidth}{0.4pt}\\[0.5cm]

\textbf{Методика обробки результатів.} Для кожного з 23 тверджень обчислюються: середнє значення ($M$), стандартне відхилення ($SD$), мінімальне та максимальне значення. Узгодженість думок експертів оцінюється за коефіцієнтом конкордації Кендалла ($W$). Прийнятним вважається $W \geq 0{,}50$ при $p < 0{,}05$.

Інтерпретація середнього значення:
\begin{itemize}
    \item $M \geq 4{,}0$~--- твердження схвалено;
    \item $3{,}5 \leq M < 4{,}0$~--- твердження потребує уточнення;
    \item $M < 3{,}5$~--- твердження потребує суттєвого перегляду.
\end{itemize}

Відкриті відповіді аналізуються методом тематичного кодування для виявлення системних зауважень та рекомендацій.

\vspace{0.5cm}

\begin{flushright}
Дата заповнення: \rule{3cm}{0.4pt}\\[0.3cm]
Підпис експерта: \rule{3cm}{0.4pt}
\end{flushright}

\label{last_page}
